% GENERATED FILE. Edit canonical lessons, then run npm run generate:books.
% canonical-source-hash: fnv1a64:abab6cfbc6530ebf
% canonical-lessons: ML-C16-kollavarsham-maasangal

\chapter{The Traditional Months}
\label{ch:ml-traditional-months}

\begin{chapteropening}
\textbf{By the end of this chapter:} \emph{I can name the twelve months of the Malayalam Kollam-era calendar and say which zodiac sign each one carries.}

Recite the twelve months from Chiṅṅam, tie Chiṅṅam to Leo, and contrast Malayalam's zodiac naming with Tamil's nakshatras.
\end{chapteropening}

\section[\ml{ചിങ്ങം} \ml{കന്നി} \ml{തുലാം} \ml{വൃശ്ചികം} \ml{ധനു} \ml{മകരം} \ml{കുംഭം} \ml{മീനം} \ml{മേടം} \ml{ഇടവം} \ml{മിഥുനം} \ml{കർക്കടകം}]{\ml{ചിങ്ങം} to \ml{കർക്കടകം} — Kerala's own solar calendar, named by the zodiac}

\label{lesson:ML-C16-kollavarsham-maasangal}



\begin{quote}
Malayalam, like Tamil, has its \textbf{own} solar calendar — but even here, the two closest cousins name their months by two \textbf{different} logics.
\end{quote}

\subsection*{What to know first}
\textbf{\ml{ചിങ്ങം}, \ml{കന്നി}, \ml{തുലാം}, \ml{വൃശ്ചികം}, \ml{ധനു}, \ml{മകരം}, \ml{കുംഭം}, \ml{മീനം}, \ml{മേടം}, \ml{ഇടവം}, \ml{മിഥുനം}, \ml{കർക്കടകം}} (\emph{Chiṅṅam, Kanni, Tulām, Vṛścikam, Dhanu, Makaram, Kumbham, Mīnam, Mēḍam, Iḍavam, Mithunam, Karkiḍakam}).

This is the \textbf{Kollam Era} calendar (also called the Malayalam calendar), Kerala's own solar civil calendar — genuinely separate from the Gregorian, the pan-Indian Hindu lunisolar calendar, \textbf{and} Tamil's solar calendar.

\begin{culture}
Both Tamil's and Malayalam's calendars are \textbf{solar}, and both are independent of the Hindu lunisolar system — but they name their months differently. Tamil ties each month to a \textbf{nakshatra} (a lunar-mansion star cluster). Malayalam instead names each month directly for the \textbf{zodiac sign} (rāśi) the sun is passing through: \textbf{\ml{ചിങ്ങം}} (\emph{Chiṅṅam}) = Leo, \textbf{\ml{കന്നി}} (\emph{Kanni}) = Virgo, \textbf{\ml{തുലാം}} (\emph{Tulām}) = Libra, and so on straight through the zodiac. Same underlying idea — track the sun's yearly path against the stars — but two Dravidian languages chose two different sets of star-markers to name it by.

\textbf{\ml{ചിങ്ങം} 1} (the 1st of Chiṅṅam) marks the \textbf{Malayalam New Year} and the start of the \textbf{Onam} festival season — Kerala's own solar calendar, running its own independent year.
\end{culture}

\subsection*{Your turn}
Say these aloud:

\begin{itemize}
\raggedright
  \item the twelve — \textquotedblleft{}Chiṅṅam, Kanni, Tulām, Vṛścikam, Dhanu, Makaram, Kumbham, Mīnam, Mēḍam, Iḍavam, Mithunam, Karkiḍakam\textquotedblright{}
  \item \textquotedblleft{}Chiṅṅam = Leo\textquotedblright{} — a zodiac sign, not a nakshatra
  \item the contrast — Tamil uses nakshatras, Malayalam uses zodiac signs
\end{itemize}

\begin{tcolorbox}[breakable,colback=teal!4,colframe=teal!35!black,title={Before you move on}]
Is Malayalam's calendar solar or lunisolar? (\textbf{Solar} — like Tamil's, and separate from the Hindu lunisolar system.) How does it differ from Tamil's in naming logic? (\textbf{Zodiac signs} — Chiṅṅam=Leo, Kanni=Virgo — rather than Tamil's nakshatra-based names.) What does Chiṅṅam 1 mark? (\textbf{Malayalam New Year}, the start of the Onam season.)
\end{tcolorbox}
